% GENERATED FILE. Edit canonical lessons, then run npm run generate:books.
% canonical-source-hash: fnv1a64:e0a71c0f3356f856
% canonical-lessons: BN-C182-chhera, BN-C182-bhaj-kora, BN-C182-joro-kora, BN-C182-bhag-kora, BN-C182-jog-kora

\chapter{To tear, To fold, To gather, To share, To add}
\label{ch:bn-to-tear-to-fold-to-gather-to-share-to-add}
\hlchaptermodality{182}

\begin{chapteropening}
\textbf{By the end of this chapter:} \emph{I can say to tear, to fold, to gather, to share and to add.}

Complete the last lesson of chapter 182: I can say to tear, to fold, to gather, to share and to add.
\end{chapteropening}

\section[\texorpdfstring{\bn{ছেঁড়া} (chhẽṛā)}{ছেঁড়া (chhẽṛā)}]{\bn{ছেঁড়া} (chhẽṛā) — to tear}

\label{lesson:BN-C182-chhera}



\begin{quote}
Before the new one: say the Bengali for to memorise, then the Bengali for to explain.
\end{quote}

\subsection*{You'll want to know: \bn{ছেঁড়া}}
\textbf{\bn{ছেঁড়া}} — \emph{chhẽṛā} — \textquotedblleft{}to tear\textquotedblright{}.

\begin{grammarlens}[title={What you've built}]
One more word for this chapter.
\end{grammarlens}

\subsection*{Your turn}
\begin{itemize}
\raggedright
  \item \emph{Say it:} \emph{chhẽṛā}
  \item \emph{Say it:} \emph{chhẽṛā}, once more
  \item \emph{Say it:} say \emph{chhẽṛā}
  \item \emph{Recall:} say the Bengali for to memorise, then the Bengali for to explain, then say \emph{chhẽṛā} again
\end{itemize}

\begin{tcolorbox}[breakable,colback=teal!4,colframe=teal!35!black,title={Before you move on}]
What does \bn{ছেঁড়া} mean? (\textquotedblleft{}To tear\textquotedblright{}.) Say it once more.
\end{tcolorbox}
\section[\texorpdfstring{\bn{ভাঁজ} \bn{করা} (bhā̃j kôrā)}{ভাঁজ করা (bhā̃j kôrā)}]{\bn{ভাঁজ} \bn{করা} (bhā̃j kôrā) — to fold}

\label{lesson:BN-C182-bhaj-kora}



\begin{quote}
Before the new one: say the Bengali for to explain, then the Bengali for to tear.
\end{quote}

\subsection*{You'll want to know: \bn{ভাঁজ} \bn{করা}}
\textbf{\bn{ভাঁজ} \bn{করা}} — \emph{bhā̃j kôrā} — \textquotedblleft{}to fold\textquotedblright{}.

\begin{grammarlens}[title={What you've built}]
One more word for this chapter.
\end{grammarlens}

\subsection*{Your turn}
\begin{itemize}
\raggedright
  \item \emph{Say it:} \emph{bhā̃j kôrā}
  \item \emph{Say it:} \emph{bhā̃j kôrā}, once more
  \item \emph{Say it:} say \emph{bhā̃j kôrā}
  \item \emph{Recall:} say the Bengali for to explain, then the Bengali for to tear, then say \emph{bhā̃j kôrā} again
\end{itemize}

\begin{tcolorbox}[breakable,colback=teal!4,colframe=teal!35!black,title={Before you move on}]
What does \bn{ভাঁজ} \bn{করা} mean? (\textquotedblleft{}To fold\textquotedblright{}.) Say it once more.
\end{tcolorbox}
\section[\texorpdfstring{\bn{জড়ো} \bn{করা} (jôṛo kôrā)}{জড়ো করা (jôṛo kôrā)}]{\bn{জড়ো} \bn{করা} (jôṛo kôrā) — to gather}

\label{lesson:BN-C182-joro-kora}



\begin{quote}
Before the new one: say the Bengali for to tear, then the Bengali for to fold.
\end{quote}

\subsection*{You'll want to know: \bn{জড়ো} \bn{করা}}
\textbf{\bn{জড়ো} \bn{করা}} — \emph{jôṛo kôrā} — \textquotedblleft{}to gather\textquotedblright{}.

\begin{grammarlens}[title={What you've built}]
One more word for this chapter.
\end{grammarlens}

\subsection*{Your turn}
\begin{itemize}
\raggedright
  \item \emph{Say it:} \emph{jôṛo kôrā}
  \item \emph{Say it:} \emph{jôṛo kôrā}, once more
  \item \emph{Say it:} say \emph{jôṛo kôrā}
  \item \emph{Recall:} say the Bengali for to tear, then the Bengali for to fold, then say \emph{jôṛo kôrā} again
\end{itemize}

\begin{tcolorbox}[breakable,colback=teal!4,colframe=teal!35!black,title={Before you move on}]
What does \bn{জড়ো} \bn{করা} mean? (\textquotedblleft{}To gather\textquotedblright{}.) Say it once more.
\end{tcolorbox}
\section[\texorpdfstring{\bn{ভাগ} \bn{করা} (bhāg kôrā)}{ভাগ করা (bhāg kôrā)}]{\bn{ভাগ} \bn{করা} (bhāg kôrā) — to share, to divide}

\label{lesson:BN-C182-bhag-kora}



\begin{quote}
Before the new one: say the Bengali for to fold, then the Bengali for to gather.
\end{quote}

\subsection*{You'll want to know: \bn{ভাগ} \bn{করা}}
\textbf{\bn{ভাগ} \bn{করা}} — \emph{bhāg kôrā} — \textquotedblleft{}to share, to divide\textquotedblright{}.

\begin{grammarlens}[title={What you've built}]
One more word for this chapter.
\end{grammarlens}

\subsection*{Your turn}
\begin{itemize}
\raggedright
  \item \emph{Say it:} \emph{bhāg kôrā}
  \item \emph{Say it:} \emph{bhāg kôrā}, once more
  \item \emph{Say it:} say \emph{bhāg kôrā}
  \item \emph{Recall:} say the Bengali for to fold, then the Bengali for to gather, then say \emph{bhāg kôrā} again
\end{itemize}

\begin{tcolorbox}[breakable,colback=teal!4,colframe=teal!35!black,title={Before you move on}]
What does \bn{ভাগ} \bn{করা} mean? (\textquotedblleft{}To share, to divide\textquotedblright{}.) Say it once more.
\end{tcolorbox}
\section[\texorpdfstring{\bn{যোগ} \bn{করা} (jog kôrā)}{যোগ করা (jog kôrā)}]{\bn{যোগ} \bn{করা} (jog kôrā) — to add}

\label{lesson:BN-C182-jog-kora}



\begin{quote}
Before the new one: say the Bengali for to gather, then the Bengali for to share.
\end{quote}

\subsection*{You'll want to know: \bn{যোগ} \bn{করা}}
\textbf{\bn{যোগ} \bn{করা}} — \emph{jog kôrā} — \textquotedblleft{}to add\textquotedblright{}.

\begin{grammarlens}[title={What you've built}]
One more word for this chapter.
\end{grammarlens}

\subsection*{Your turn}
\begin{itemize}
\raggedright
  \item \emph{Say it:} \emph{jog kôrā}
  \item \emph{Say it:} \emph{jog kôrā}, once more
  \item \emph{Say it:} say \emph{jog kôrā}
  \item \emph{Recall:} say the Bengali for to gather, then the Bengali for to share, then say \emph{jog kôrā} again
\end{itemize}

\begin{tcolorbox}[breakable,colback=teal!4,colframe=teal!35!black,title={Before you move on}]
What does \bn{যোগ} \bn{করা} mean? (\textquotedblleft{}To add\textquotedblright{}.) Say it once more.
\end{tcolorbox}
